%%% ==========================================================================
%%%  macros-math.tex —— 数学讲义专用宏
%%%  仅收录跨章节反复使用的记号；一次性记号请就地定义在所在章内。
%%% ==========================================================================

%% ---- 数系与集合 ----
\newcommand{\N}{\mathbb{N}}
\newcommand{\Z}{\mathbb{Z}}
\newcommand{\Q}{\mathbb{Q}}
\newcommand{\R}{\mathbb{R}}
\newcommand{\C}{\mathbb{C}}
\newcommand{\F}{\mathbb{F}}
\newcommand{\set}[1]{\left\{#1\right\}}
\newcommand{\setst}[2]{\left\{\,#1 \;\middle|\; #2\,\right\}}   % {x | P(x)}

%% ---- 常用算子 ----
\DeclareMathOperator{\im}{Im}                  % 像
\DeclareMathOperator{\rank}{rank}
\DeclareMathOperator{\tr}{tr}
\DeclareMathOperator{\spn}{span}
\DeclareMathOperator{\diag}{diag}
\DeclareMathOperator{\sgn}{sgn}
\DeclareMathOperator{\supp}{supp}
\DeclareMathOperator*{\esssup}{ess\,sup}
\renewcommand{\ker}{\operatorname{Ker}}

%% ---- 括号与范数（mathtools 可伸缩版本）----
\DeclarePairedDelimiter{\abs}{\lvert}{\rvert}
\DeclarePairedDelimiter{\norm}{\lVert}{\rVert}
\DeclarePairedDelimiter{\inner}{\langle}{\rangle}
\DeclarePairedDelimiter{\floor}{\lfloor}{\rfloor}
\DeclarePairedDelimiter{\ceil}{\lceil}{\rceil}

%% ---- 微积分 ----
\newcommand{\dif}{\mathop{}\!\mathrm{d}}       % 正体微分号（国标要求）
\newcommand{\pd}[2]{\frac{\partial #1}{\partial #2}}
\newcommand{\dd}[2]{\frac{\dif #1}{\dif #2}}
\newcommand{\eval}[1]{\left.#1\right|}
\newcommand{\eu}{\mathrm{e}}                   % 自然常数（国标正体）
\newcommand{\iu}{\mathrm{i}}                   % 虚数单位（国标正体）

%% ---- 向量与矩阵 ----
\newcommand{\vect}[1]{\boldsymbol{#1}}
\newcommand{\mat}[1]{\boldsymbol{#1}}
\newcommand{\transpose}{^{\mathsf{T}}}
\newcommand{\inv}{^{-1}}

%% ---- 数列 ----
\newcommand{\seq}[1]{\{#1\}}                     % 数列 \seq{a_n} 排作 {a_n}

%% ---- 收敛与映射箭头 ----
\newcommand{\mapsdef}{\colon}                  % f\mapsdef X \to Y
\newcommand{\converge}[1]{\xrightarrow{\;#1\;}}
